\chapter{INTRODUCCIÓN}

La industria gastronómica colombiana enfrenta una contradicción estructural: a pesar de 
constituir uno de los sectores productivos más dinámicos del país, con aproximadamente 
167.000 establecimientos registrados~\cite{acodres_estadisticas_2021} y un aporte del 
3,9\% al Producto Interno Bruto nacional equivalente a 56,7 billones de pesos en 
2022~\cite{andi_pib_2024}, la gran mayoría de sus operaciones cotidianas de gestión de 
pedidos se sostienen sobre un procedimiento manual: el mesero anota la orden del cliente 
en papel y la entrega físicamente a la cocina. Este método, históricamente aceptado por 
su bajo costo operativo inmediato, genera tres categorías de problemas que comprometen la 
viabilidad de los establecimientos de pequeña y mediana escala: ineficiencia operativa, 
riesgo sanitario e inseguridad de la información, todos susceptibles de ser mitigados 
mediante la adopción de tecnología digital adecuada a las condiciones reales del sector.

La dimensión operativa del problema está respaldada por evidencia cuantitativa concreta. 
Los sistemas basados en papel presentan tasas de error en pedidos de hasta el 8\%, frente al 2\% registrado en sistemas digitales, siendo la causa principal la mala 
interpretación de la caligrafía y el uso de abreviaturas en entornos de trabajo 
acelerado~\cite{Daniel2024, LSRetail2025}. La Asociación Nacional de Restaurantes de 
los Estados Unidos estima que este tipo de fallos puede traducirse en pérdidas equivalentes 
al 4\% de las ventas totales del establecimiento~\cite{Khan2023}, dado que 
aproximadamente el 60\% de los errores resultan en reemplazos gratuitos que el negocio 
debe asumir~\cite{Checkmate2025}. Esta presión económica adquiere una dimensión crítica 
en el contexto colombiano, donde la Asociación Colombiana de la Industria Gastronómica 
(ACODRES) estima que cuatro de cada diez restaurantes cierran durante sus primeros seis 
meses de operación~\cite{acodres_estadisticas_2021}. La imposibilidad de actualizar 
órdenes en tiempo real limita además la personalización del servicio y compromete la 
fidelización del cliente~\cite{Ser2024Estrategia}. A pesar de que las mejoras tecnológicas 
son conocidas, la implementación de sistemas digitales comerciales requiere una inversión 
inicial elevada en hardware, software y capacitación, sumada a costos de mantenimiento 
continuos que resultan poco accesibles para establecimientos con recursos 
limitados~\cite{marcelo2024desafios}.

La dimensión sanitaria, menos visible pero igualmente documentada, está asociada al papel 
como vehículo de transmisión de microorganismos. Stephens et al.~\cite{Stephens2019} 
demostraron en una revisión sistemática que los fómites constituyen un mecanismo de 
transmisión relevante para numerosos agentes infecciosos, y que las características físicas 
del soporte determinan directamente la tasa de supervivencia microbiana. Kirchner 
et al.~\cite{Kirchner2021} confirmaron que las superficies atípicas en entornos de 
servicio de alimentos, es decir, aquellas que quedan fuera de los protocolos de limpieza 
convencionales, actúan como vectores activos de contaminación cruzada. La comanda de papel 
recorre un circuito de múltiples manos (mesero, cocinero y cajero) sin protocolos de 
desinfección intermedios, y Hübner et al.~\cite{Huebner2011} demostraron experimentalmente 
que bacterias como \textit{Pseudomonas aeruginosa} y \textit{Enterococcus hirae} 
sobreviven en papel estándar durante al menos siete días y son recuperables mediante el 
ciclo de contacto mano--papel--mano. En el contexto colombiano, un análisis microbiológico 
en restaurantes formales e informales de Cali confirmó contaminación cruzada con altos 
porcentajes de bacterias mesófilas aerobias en superficies de contacto con 
alimentos~\cite{caro-hernandez_analisis_2020}, situación que contraviene directamente la 
Resolución 2674 de 2013 del Ministerio de Salud y Protección Social, que exige minimizar 
cualquier foco de contaminación que ponga en riesgo la inocuidad 
alimentaria~\cite{minsalud_resolucion2674_2013}.

La tercera dimensión del problema radica en la seguridad de la información. Las soluciones 
comerciales de gestión de comandas disponibles en el mercado colombiano (entre ellas 
GastroPOS~\cite{gastropos}, Loggro~\cite{loggro_planes}, Vendty~\cite{vendty} y 
Toteat~\cite{toteat_planes}) operan mayoritariamente bajo esquemas de suscripción mensual 
en la nube, sin contemplar de forma explícita mecanismos robustos de encriptación ni 
garantías sobre la custodia soberana de los datos del establecimiento. Esta estructura 
genera desconfianza en el operador, pues el control de información sensible (historiales 
de ventas, inventarios y preferencias de clientes) queda en manos de terceros. Para las 
PYME gastronómicas, que según la Encuesta de Micronegocios del DANE (EMICRON, primer 
trimestre de 2024) enfrentan limitaciones estructurales de inversión 
tecnológica~\cite{DANE2024}, esta barrera económica e informática consolida la resistencia 
a la digitalización y deja al sector en una situación de vulnerabilidad operativa y 
sanitaria sostenida.

Investigadores y desarrolladores de distintas latitudes han abordado problemas 
parcialmente similares al planteado en este trabajo. En el ámbito académico colombiano, 
Corredor Siratá~\cite{Documento2022} propuso un sistema POS web para la gestión de 
pedidos, evidenciando la necesidad de digitalización en el sector pero sin incorporar 
mecanismos de seguridad multicapa ni operación en entorno local. Desde España, Curieses 
Fernández~\cite{Curieses2023} y Bernat Elorza~\cite{Bernat2023} desarrollaron sistemas 
de gestión de comandas mediante aplicaciones web y dispositivos móviles, priorizando la 
usabilidad pero sin abordar la restricción de conectividad ni la protección de datos 
sensibles. En Latinoamérica, Escobar Vallejos~\cite{Escobar2023} y Allain 
Manrique~\cite{Allain2020} construyeron prototipos funcionales para establecimientos 
específicos, mientras que De Argila Lorente~\cite{deArgila2024} analizó la evolución 
de la gestión de comandas en tiempo real como parte de la transformación digital de la 
restauración española. En el plano técnico y operativo, Katagri y 
Balavalad~\cite{Katagri2025} documentaron reducciones del 40\% en tiempos de 
procesamiento de órdenes mediante automatización digital, y Sultana et 
al.~\cite{Sultana2024} propusieron un modelo basado en IoT para mejorar las operaciones 
de servicio de alimentos en países en desarrollo. Hummel y Murphy~\cite{EmilyHummel} 
demostraron mediante la técnica de planos de servicio que la estandarización de los 
canales de comunicación interna mejora directamente la eficiencia operativa, y De Vries 
et al.~\cite{DEVRIES201859} confirmaron que los sistemas digitales de pedido y pago 
impactan la retención de clientes y la frecuencia de visitas. Liyanage 
et al.~\cite{Liyanage2018} abrieron el precedente técnico de integración de tecnología 
NFC bajo el estándar ISO/IEC~14443 en sistemas de gestión de restaurantes. Sin embargo, 
ninguno de estos antecedentes aborda de forma integrada la restricción económica del 
hardware, la operación sin conectividad a internet, la encriptación de datos y la 
autenticación física como unidad funcional orientada específicamente al segmento PYME 
colombiano.

La pregunta que orienta esta investigación es: \textit{¿cómo desarrollar una aplicación 
	web de bajo costo para la gestión de comandas en restaurantes que garantice la seguridad 
	y confidencialidad de la información mediante encriptación y llaves físicas de acceso?} 
El objeto de estudio es el diseño, implementación y evaluación de \textbf{QuickTable}, 
un sistema web de gestión integral de comandas que digitaliza el ciclo operativo completo 
del restaurante (toma de pedidos, gestión en cocina y facturación en caja), incorporando 
encriptación de base de datos y autenticación de doble factor (2FA) basada en hardware 
NFC, operable sobre infraestructura de bajo costo en red de intranet local sin dependencia 
de conectividad externa. La hipótesis central del proyecto, planteada en el anteproyecto, 
sostiene que es posible construir una solución tecnológica accesible, segura y funcional 
para el sector gastronómico colombiano sin recurrir a servicios en la nube ni a hardware 
especializado de alto costo, siempre que se combinen adecuadamente tecnologías de software 
maduras con componentes electrónicos de bajo costo orientados a la seguridad por hardware.

El desarrollo de este proyecto se justifica por el impacto que genera en múltiples planos. 
En el plano operativo y económico, los restaurantes PYME que adopten la solución reducirán 
errores de pedido, acortarán tiempos de procesamiento de órdenes y obtendrán trazabilidad 
histórica de sus ventas con un costo de implementación sustancialmente inferior al de las 
alternativas comerciales existentes; los beneficiarios directos son los propietarios y 
administradores de establecimientos gastronómicos de pequeña y mediana escala, así como su 
personal operativo. En el plano sanitario, la eliminación del papel como medio de 
comunicación interna rompe activamente la cadena de transmisión de microorganismos entre 
áreas, constituyendo una medida preventiva que protege tanto al personal del 
establecimiento como a los clientes~\cite{minsalud_resolucion2674_2013}. En el plano 
informático, el modelo de despliegue local garantiza que el restaurante mantenga la 
custodia absoluta de su información financiera y operativa, y la encriptación de la base de 
datos junto con el módulo de autenticación física elevan el nivel de protección a un 
estándar que las soluciones comerciales del mercado colombiano no contemplan de forma 
integrada. Adicionalmente, el proyecto realiza un aporte metodológico al demostrar que los 
principios de la Ingeniería Electrónica (comunicaciones inalámbricas, sistemas embebidos y 
arquitecturas de seguridad) son aplicables a la solución de problemas reales del sector 
productivo nacional, diferenciando esta propuesta de los desarrollos puramente 
informáticos previos.

Para resolver el problema planteado se adoptó un enfoque de ingeniería orientado a la 
construcción e implementación real del sistema siguiendo una metodología en cascada. La 
aplicación web fue desarrollada sobre el framework .NET~8 con el modelo de páginas Razor 
Pages y Entity Framework~6~\cite{ef6_docs} como capa de acceso a datos sobre un motor 
Microsoft SQL Server, con una interfaz de usuario estructurada sobre la plantilla 
AdminLTE~\cite{BALSAM2025112186}. El módulo de seguridad por hardware se implementó 
mediante una Raspberry Pi~3 equipada con un lector RFID-RC522 operando a 13,56\,MHz 
conforme al estándar ISO/IEC~14443~\cite{Liyanage2018}, gestionada por un cliente Python 
con interfaz gráfica desarrollada en Tkinter. La validación del sistema se realizó mediante 
pruebas unitarias con MSTest, análisis de vulnerabilidades con OWASP~ZAP~2.17.0, pruebas 
automatizadas de concurrencia y una implementación piloto en entorno real sobre 
infraestructura deliberadamente limitada, con el fin de verificar la viabilidad del sistema 
en las condiciones reales del segmento objetivo.